\section{Probes on the reference cells}\label{sec:reference-cells}

The sampling and placement arguments are those of
\cite[\S19]{QED3Master}. The external-source identity uses a finitely supported bond cochain.
The continuum cumulant uses the test $f=dd^*\beta$. We identify the
discrete source with a nearby smooth test and place it strictly inside
one reference cell. Copies in distinct cells then have orthogonal
curvatures, which makes the normal-ordering factor separate over cells.

\subsection{The discrete source and the smooth test}

Fix a real nonzero $\beta\in C_c^\infty(\R^3;\Lambda^2)$.
Use the unbuffered support-cube side $D_\beta$ from the introduction,
and write
$D_\beta^{\mathrm{buf}}=D_\beta+4$. Choose a cube
$Q_\beta=z_\beta+[-D_\beta^{\mathrm{buf}}/2,
D_\beta^{\mathrm{buf}}/2]^3$ containing
$\supp\beta+[-2,2]^3$. Use the centered differences
$\delta_{\eta,\mu}$ and cube average $J_\eta$ defined in the introduction.
Put
\begin{align}
 a'_{\eta,\nu}&=-\sum_{\mu\ne\nu}
                          \delta_{\eta,\mu}\beta_{\mu\nu},\notag\\
 f^\eta_{\mu\nu}&=\delta_{\eta,\mu}a'_{\eta,\nu}
                         -\delta_{\eta,\nu}a'_{\eta,\mu}.
                                      \label{eq:cells-smooth-test}
\end{align}
All these functions are smooth and real. Since each centred
difference enlarges support by at most $\eta/2$ in its direction,
\[
 \supp f^\eta,\ \supp J_\eta a'_\eta
                   \subset\supp\beta+[-\eta,\eta]^3.
\]
In particular their supports lie in $Q_\beta$ for $0<\eta\le1$.

On the physical lattice $\eta\Z^3$ set
\[
 b_{\eta,\mu\nu}(c)
       =(J_\eta\beta_{\mu\nu})
                       (c+\eta(e_\mu+e_\nu)/2),\qquad
 a_\eta^o=d_\eta^*b_\eta.
\]
Here $d_\eta$ is the mesh-normalized forward cochain differential
and $d_\eta^*$ its adjoint for the pairing $\eta^3\sum$.
The adjoint formula gives
\begin{equation}\label{eq:cells-midpoint-source}
 a_{\eta,\nu}^o(c)=(J_\eta a'_{\eta,\nu})(c+\eta e_\nu/2).
\end{equation}
Transport to the unit lattice and define
\[
 \hat{\mathfrak a}_\eta^o(x,x+e_\nu)
                    =\eta^{1/2}a_{\eta,\nu}^o(\eta x).
\]
Its coboundary uses unscaled unit-lattice differences. The
midpoint identities yield
\begin{equation}\label{eq:cells-source-curvature}
 d\hat{\mathfrak a}_\eta^o(\sigma)
                        =f^\eta_{[\eta]}(\sigma).
\end{equation}
Indeed, the bond values in the curl are evaluated at the two
points displaced by $\pm\eta e_\mu/2$ from the plaquette centre.
Their difference is $\eta\delta_{\eta,\mu}J_\eta a'_{\eta,\nu}$.
Commutation of the average and differences then gives
$\eta^{3/2}J_\eta f^\eta$ at that centre, which is the right
side of \eqref{eq:cells-source-curvature}.

For a compactly supported smooth two-form $g$, let $g_{[\eta]}$
be the covering plaquette cochain with value $\eta^{3/2}J_\eta g$
at the plaquette center. On a torus $C(n)$ of physical side
$R=n\eta$, periodize this cochain by summing its translates by
$n\Z^3$. Write $F^\sharp$ for the periodic componentwise step
field whose value on the physical cube centered at a plaquette
is $\eta^{-3/2}F^{\rm u}$. Then $F^\sharp(g)$ is the finite-torus
pairing of the periodized $g_{[\eta]}$ with $F^{\rm u}$.
Periodizing the source likewise gives the exact observable identity
\begin{equation}\label{eq:cells-observable}
 \langle d\hat{\mathfrak a}_\eta,F^{\mathrm u}\rangle
                            =F^\sharp(f^\eta).
\end{equation}
The source is transverse: $a_\eta^o=d_\eta^*b_\eta$ is
orthogonal to every closed bond cochain, by summation by
parts. This applies to both the covering cochain and its
periodization.

The continuum approximation has the uniform bound
\begin{equation}\label{eq:cells-test-error}
 \|f^\eta-f\|_{H^s}\le\eta^2\|\beta\|_{H^{s+4}}
                         \qquad(s\in\R,\ 0<\eta\le1).
\end{equation}
To see this, the difference symbol is
$is_\mu(p)$ with $s_\mu(p)=2\eta^{-1}\sin(\eta p_\mu/2)$.
One has $|s_\mu|\le|p_\mu|$ and
$|s_\mu-p_\mu|\le\eta^2|p_\mu|^3/24$, so
\[
 |s_\mu s_\nu-p_\mu p_\nu|
                               \le\eta^2|p|^4/12.
\]
Each component of $f^\eta-f$ contains four such terms.
Cauchy--Schwarz on these terms, followed by summation over
the three two-form components, gives
$|\widehat{f^\eta-f}|^2\le(\eta^2/3)^2|p|^8|\hat\beta|^2$.
Plancherel's theorem proves \eqref{eq:cells-test-error}.

\subsection{Reference geometry and placement}

Let $T^*\ge1$ be the reference exponent furnished by the
construction, and put
\[
 R^*=hL^{T^*},\qquad n_*=R^*/\eta,\qquad
 p_{\rm cell}=n_*/L,\qquad w=p_{\rm cell}/2,\qquad \ell_c=\eta w=R^*/(2L).
\]
Since $h$ and $\eta^{-1}$ are powers of the odd integer $L$,
$p_{\rm cell}$ is an odd integer, $w\in\Z+1/2$, and the open unit-lattice
reference cells are $Q_c^\circ=w(c+(0,1)^3)$.
The reference torus $C(n_*)$ has $2L$ cells per side and
$N_*=8L^3$ in all.

The probe fits the reference cells if
\begin{equation}\label{eq:cells-fit}
 D_\beta^{\mathrm{buf}}+3\eta\le\ell_c.
\end{equation}
A placement is a vector $v\in\eta\Z^3$ satisfying
\[
 |v+z_\beta-(\ell_c/2)(1,1,1)|_\infty\le\eta/2.
\]
Such a vector always exists by coordinatewise rounding.
It translates $Q_\beta$ into $[\eta,\ell_c-\eta]^3$:
the displacement of its centre is at most $\eta/2$, and
the half-side is at most $(\ell_c-3\eta)/2$.
Hence
\begin{equation}\label{eq:cells-probe-margin}
 \supp(\tau_v\beta)
                     \subset[2+\eta,\ell_c-2-\eta]^3.
\end{equation}
Form the source $\hat{\mathfrak a}^o$ with this translated
probe.

Every endpoint of a bond in the source support is within
$3\eta/2$ in sup distance of $\supp(\tau_v\beta)$.
By \eqref{eq:cells-probe-margin}, its lattice coordinates
therefore lie in
$[2/\eta-1/2,w-2/\eta+1/2]$.
If a plaquette has nonzero $d\hat{\mathfrak a}^o$, at least
one of its boundary bonds belongs to that support. Its
other two corners are one unit away. Thus the whole
plaquette lies in the box
\begin{equation}\label{eq:cells-support-box}
 B_0=(w/2)(1,1,1)+[-r_0,r_0]^3,\qquad
 r_0=w/2-2/\eta+3/2.
\end{equation}
Condition \eqref{eq:cells-fit} gives $r_0\ge3$.
The distance of this box from each bounding wall is
$2/\eta-3/2\ge1/2$, so $B_0\subset Q_0^\circ$.
Its reflected copies lie strictly inside the corresponding
cells. This is the support margin needed by the local
source identity.

\subsection{Orthogonal sources and translated observables}

For a cell index $c$, define the wall map $g_c$ coordinatewise:
if $c_\mu=2i$, it sends $x_\mu$ to $x_\mu+2iw$; if
$c_\mu=2i+1$, it sends $x_\mu$ to $(2i+2)w-x_\mu$.
It carries $Q_0^\circ$ to $Q_c^\circ$ and acts on cochains by
transporting oriented bonds, with a sign upon orientation reversal. Put
$\hat{\mathfrak a}_c^o=(g_c)_*\hat{\mathfrak a}^o$ on the
covering lattice and periodize it on $C(n)$, where $p_{\rm cell}$ divides
$n$ and $n\ge n_*$, to obtain $\hat{\mathfrak a}_c$.
The period lattice identifies cell indices modulo
$(2n/p_{\rm cell})\Z^3$. The coboundary commutes with these actions
and periodization, and its support lies in $B_c=g_cB_0$.
For cell representatives on this torus,
\begin{equation}\label{eq:cells-orthogonality}
 \langle d\hat{\mathfrak a}_c,d\hat{\mathfrak a}_{c'}\rangle_n
             =v_\eta\,\mathbf1_{\{c=c'\}},\qquad
 v_\eta=\|d\hat{\mathfrak a}^o\|^2.
\end{equation}
The covering norm is the unweighted $\ell^2$ norm, and the
torus pairing is the counting pairing.

To prove the identity, unfold one factor:
\[
 \langle d\hat{\mathfrak a}_c,d\hat{\mathfrak a}_{c'}\rangle_n
  =\sum_{\sigma}\sum_{z\in\Z^3}
       d\hat{\mathfrak a}_c^o(\sigma)
       d\hat{\mathfrak a}_{c'}^o(\sigma+nz).
\]
A nonzero summand has $\sigma\subset Q_c^\circ$ and
$\sigma+nz\subset Q_{c'}^\circ$. The open cells are disjoint,
so necessarily $c'=c+(2n/p_{\rm cell})z$. For distinct torus cells
there is no such $z$. For the same cell exactly one translate
can contribute; its summands are the squares of
$d\hat{\mathfrak a}_c^o$. The wall action preserves the
counting norm, which proves \eqref{eq:cells-orthogonality}.
Changing the placement by a vector in $\eta\Z^3$ translates
the covering source by an integer vector. Thus $v_\eta$
depends only on $\beta$ and $\eta$, and is independent of
the torus and the placement.

For a closed cochain $u$ define
$X_c(u)=\langle d\hat{\mathfrak a}_c,u\rangle_n$.
Every polynomial in this observable is read by one and the
same local integrand on all these tori. On the cover its
local expression is
\[
 \widehat A\longmapsto
 P\left(\sum_{\sigma\subset B_c}
             d\hat{\mathfrak a}_c^o(\sigma)(d\widehat A)(\sigma)\right).
\]
This is a scalar gauge-invariant polynomial. Cauchy--Schwarz
bounds its argument by
\[
 \sqrt{v_\eta}\left(\sum_{\sigma\subset B_c}
                      |d\widehat A(\sigma)|^2\right)^{1/2}.
\]
The polynomial is therefore bounded by
\[
 C_{\mathrm{pol}}\exp\left(\frac14\sum_{\sigma\subset B_c}
                                      |d\widehat A(\sigma)|^2\right)
\]
for a finite constant $C_{\mathrm{pol}}$ depending on the polynomial
and the covering source. This is the local Gaussian-growth condition used for the
observable class $\mathcal A(Q_c^\circ)$ in \cite[\S4.1, (46)]{ChamberLattice}.

The local boxes embed in the torus. Indeed, writing
$p_{\rm cell}=2m_{\rm cell}+1$, condition \eqref{eq:cells-fit} gives $m_{\rm cell}\ge7$,
and
\[
 2r_0+2=w-4/\eta+5\le w+1<p_{\rm cell}\le n.
\]
For each $c$, the translates of $B_c+[-1,1]^3$ by
$nz$, $z\in\Z^3$, are consequently disjoint.
On a local lift, $d\widehat A=F^{\mathrm u}$; unfolding
the finite cochain sum reads the polynomial as $P(X_c)$.
For a real polynomial its reflected copy from the base
cell is exactly $P(X_c)$, since conjugation fixes its
coefficients. This applies to the centered Wick polynomials.

The variance in \eqref{eq:cells-orthogonality} approaches
the continuum normal-ordering variance:
\begin{equation}\label{eq:cells-variance-limit}
 \left|\sqrt{v_\eta}-\|f\|_2\right|
 \le21\eta|\beta|_{C^3}(D_\beta^{\mathrm{buf}})^{3/2}
                      +\eta^2\|\beta\|_{H^4}.
\end{equation}
Let $\Pi_\eta f^\eta$ be the componentwise cell-average step
function on a period cube. By \eqref{eq:cells-source-curvature},
its $L^2$ norm is $\sqrt{v_\eta}$.
Each first derivative of $f^\eta_{\mu\nu}$ is bounded
by $4|\beta|_{C^3}$, from the four centred-difference
terms in \eqref{eq:cells-smooth-test} and the identity
$\delta_{\eta,\mu}g=\int_{-1/2}^{1/2}
(\partial_\mu g)(\,\cdot+t\eta e_\mu)\,\dd t$.
The pointwise error of the cell average is at most
$12\eta|\beta|_{C^3}$. It vanishes outside the periodic
images of the buffered cube, whose measure in one period
is $(D_\beta^{\mathrm{buf}})^3$.
Summing the three component squares bounds the $L^2$
error by $21\eta|\beta|_{C^3}(D_\beta^{\mathrm{buf}})^{3/2}$.
The norm triangle inequality and
\eqref{eq:cells-test-error} prove
\eqref{eq:cells-variance-limit}.

Finally the even cell indices give pure translations.
For $y\in(R^*/L)\Z^3=p_{\rm cell}\eta\Z^3$ set
$c(y)=2y/(p_{\rm cell}\eta)\in2\Z^3$. Then $g_{c(y)}$ is the unit-lattice
translation by $y/\eta$, and
\begin{equation}\label{eq:cells-translation-observable}
 X_{c(y)}=F^\sharp(\tau_{v+y}f^\eta).
\end{equation}
Two such cells agree on $C(n)$ exactly when their physical
translations agree modulo $R\Z^3$. Any fixed finite set
of distinct translations therefore occupies distinct
cells once the volume is sufficiently large.
Lattice translation invariance removes the common placement
$v$ from the expectation of any polynomial in these
observables. Its integrability follows from Gaussian
domination. These identities enter the subsequent volume transfer of
\cite[Theorem~A]{VolumeTransfer}, which supplies the larger-volume estimates
for the translates $\tau_yf^\eta$ and their passage to component cumulants.


